\section{Introduction} \label{sec:introduction}

\acr{LLM}s have recently made impressive
advancements across a range of domains and tasks
\citep{anthropicIntroducingClaude35,googleIntroducingGemini2024,openaiOpenAIO1System}.
However, in order to put these \acr{LLM}s to use in real world applications,
\acr{LLM}s must be wrapped in
code to orchestrate them and expose tools that allow the models to take
actions. These action-taking \acr{LLM}s are referred to as agents, and the
broader system an \emph{agent system}.

These agent systems often show dramatic improvements in benchmark performance
over ``plain'' \acr{LLM}s \citep{yaoReActSynergizingReasoning2023,
zelikmanParselAlgorithmicReasoning2023, chenProgramThoughtsPrompting2023},
through combinations of prompting strategies and methods for combining
different \acr{LLM} outputs. Early examples include best-of-$N$ sampling and
simple prompting strategies such as chain of thought
\citep{kojimaLargeLanguageModels2022}. However more sophisticated schemes have
shown success in getting the desired behavior and
performance improvements from the models, for instance STaR
\citep{zelikmanSTaRBootstrappingReasoning2022}, Tree of Thoughts
\citep{yaoTreeThoughtsDeliberate2023}, Graph of Thoughts
\citep{bestaGraphThoughtsSolving2023}, \acr{LLM} Debate
\citep{duImprovingFactualityReasoning2023}, Iterative Self-Refinement
\citep{madaanSelfRefineIterativeRefinement2023}, Expert Prompting
\citep{longMultiexpertPromptingImproves2024} among many others.
The comprehensive survey of \citet{schulhoffPromptReportSystematic2024}
demonstrates the vast number of manually created strategies to date.

Recent improvements in coding agents raise the question of whether these agents themselves can autonomously modify
and improve their own code by discovering e.g.\ new prompting schemes or tools without manual design and implementation.
We argue that this style of fully self-referential meta-agent programming is
possible today and offers a sound alternative to the ad-hoc, trial-and-error
approach of hand-crafted orchestrators, which may only explore a small fraction
of the solution space.
Recent work in the Automated Design of Agentic Systems (\acr{ADAS})
\citep{huAutomatedDesignAgentic2024} uses a meta-agent to optimise agent
implementations.
However, \citet{huAutomatedDesignAgentic2024} is not \textit{self}-improving, as
there are two separate agents: the target-agent that performs the task, and the
meta-agent, which improves the target agent. A motivation for a self-improving
system is that the improvements in coding abilities may be leveraged during
subsequent improvement steps, hopefully compounding.

Our contributions are:
\vspace{-5pt}
\begin{itemize}
\item  A self-improving coding agent (\acr{SICA}) that eliminates the
distinction between meta-agent and target agent, and is capable of editing its own codebase to
improve itself with respect to its cost, speed and benchmark performance.
\item Empirical evidence that self-referential agents can effectively improve
      their own implementations; we find performance improves from 17\% to 53\% performance on a random subset of SWE-Bench Verified, even with consideration given to safety constraints and
      resource efficiency.
\item We share a our implementation of a self-improving coding agent (\acr{SICA}) with the community.  SICA is implemented in standard Python without a domain-specific language, and provides a reference agent framework for building new SICA systems, as well as those seeking to post-train LLMs on tool use and other agent tasks.
\end{itemize}

We make our code available at {\small \url{https://github.com/MaximeRobeyns/self\_improving\_coding\_agent}}.
